\section{Problem Setup}
\label{sec:setup}

What must the information structure respect for a sequential test to be
legitimate inside a closed control loop?

Let $I_t$ be on-hand inventory at the start of period $t$, $P_t$ the aggregate
in-transit total, $A_t$ our order, $R_t$ aggregate receipts, and $p$ and $h$ the
per-unit profit and holding cost. The order is coerced and capped as
$A_t = \min(\max(0, \lceil q_t \rceil), C_t)$ and dispatched with a latent
transit time, and our endpoints accumulate profit and lost sales under the audited
event order. Let $Y_t$ be the exposed observable,
uncensored demand in the headline setting, and $\Filt_t$ the sigma-algebra
generated by every observation through period $t$, taken \emph{immediately
before} $A_t$ is chosen. Two properties follow. \textbf{(A1) Actions are
predictable:} $A_t$ is $\Filt_t$-measurable, and the runner builds and freezes
the observation before calling the controller, so this holds by construction
rather than by convention. \textbf{(A2) Actions affect only the future:} $A_t$
influences $Y_{t+1}$ and later, never $Y_t$. Conditioning $Y_r$ on
$(\Filt_{r-1}, A_{r-1})$ is therefore well defined and does not condition on
anything $A_{r-1}$ could have altered retroactively, and every density below is
written that way. A zero actual lead time moves $R_t$ rather than $Y_t$, so (A2)
survives. Finally, $\tau_j$ is the proposal period of hypothesis $j$ and a
stopping time, $h_j$ the \shockspec{} frozen there, and $E_{j,t}$ its evidence
process.
